
%\needspace{4\baselineskip}
\color{uniblue}{\subsection{Защита от дуговых замыканий (ЗДЗ)}}
\color{black}

\begin{enumerate}

\item
Защита от дуговых замыканий предназначена для быстродействующего отключения или сигнализации коротких замыканий внутри ячеек КРУ/ЗРУ, сопровождающихся электрической дугой.

\item 
В устройстве защита от дуговых замыканий представлена функцией «ЗДЗ». Ввод функции в работу на этапе параметрирования устройства осуществляется программным переключателем SGF1 «Ввод\_функции». Информация о введенной функции отображается выходным сигналом «ЗДЗ: Ввод».

\item
При активном состоянии входного сигнала «ЗДЗ на сигн» действие защиты переводится на сигнал.

\item
Функция может быть выведена из работы оперативно путем активации сигнала «ОВ ЗДЗ», что характеризуется активным состоянием сигнала «ЗДЗ: Оперативный вывод» на выходе функции.

\item
Функция контролирует сигнал отключения от дуговой защиты «Срабатывание ЗДЗ». Для отстройки от ложной работы дуговой защиты необходимо подтверждение срабатывания от ИО тока и напряжения нулевой последовательности, которое контролируется положением уставки SGF3.

\item
Информация о превышении максимального из токов или напряжения нулевой последовательности выше уставок будет отображаться выходными сигналами «ЗДЗ: ИО 3U0» и «ЗДЗ: ИО Iмакс».

\item
Пусковые условия функции формируются при наличии входного дискретного сигнала «Срабатывание ЗДЗ», отсутствие неисправности ЗДЗ и выполнении следующих условий:
\begin{itemize}
    \item уставка SGF3 находится в положении «Не предусмотрено», без контроля тока и напряжения;
    \item если уставка SGF3 находится в положении «Контроль I», превышение максимального из токов выше уставки;
    \item если уставка SGF3 находится в положении «Контроль I или 3U0», превышение максимального из токов или напряжения НП выше уставок при отсутствии сигнала неисправности в цепях напряжения при введенной уставке SGF4,  предусматривающей контроль неисправности в цепях напряжения (БНН);
    \item если уставка SGF3 находится в положении «Внешний», наличие входного дискретного сигнала «Контроль тока ЗДЗ» без контроля по току и напряжению.
\end{itemize}

\item
При длительном наличии сигнала «Срабатывание ЗДЗ» дольше выдержки времени T2 формируется сигнал о наличии неисправности ЗДЗ «ЗДЗ: Неисправность». Наличие сигнала неисправности ЗДЗ блокирует пусковые условия функции при введенной уставке SGF2. Дискретный сигнал «Срабатывание ЗДЗ» поступает из логики дуговой защиты информируя о засвете волоконнооптического датчика.

\item
Алгоритм работы функции выполнен в соответствии с рисунком \ref{zdz:img1}. В таблице \ref{zdz:tbl1} приведены параметры, необходимые для настройки функции.

\vspace{1mm}
\begin{figure}[H]
\centering
\includegraphics[width=1\textwidth,height=1\textheight,keepaspectratio]{\fbpath/108. ЗДЗ/_latex/img/zdz.pdf}
\captionof{figure}{Функциональная схема логики ЗДЗ}\label{zdz:img1}
\end{figure}

\end{enumerate}
